\hypertarget{appendix-8bc4e8fc547c}{}\label{referencestart:appendix-8bc4e8fc547c}
\section{Dot products and orthogonality}\label{appendix-8bc4e8fc547c-dot-products-and-orthogonality}

\hypertarget{appendix-8bc4e8fc547c-dot-products-and-orthogonality}{}

\textbf{Curriculum depth:} Introduced · \textbf{First seen:} Phase6 · \textbf{Latest development:} Phase6

\begin{allmightroledefinition}

\subsection{Definition and mechanism}\label{appendix-8bc4e8fc547c-definition-and-mechanism}

\hypertarget{appendix-8bc4e8fc547c-definition-and-mechanism}{}

An inner product provides a scalar measurement of two vectors. In orthonormal real coordinates the dot is a sum of component products. Its geometric form is length product times angle cosine. Orthogonality means zero inner product; zero itself has no angle.

\end{allmightroledefinition}

\((3,4,0)\cdot(2,0,0)=6\), while the cosine is \(6/(5\cdot2)=0.6\). Coordinate dotting assumes the correct metric.

\subsection{Prerequisites}\label{appendix-8bc4e8fc547c-prerequisites}

\hypertarget{appendix-8bc4e8fc547c-prerequisites}{}

\begin{itemize}
\tightlist
\item
  \href{../07-Vision-XR-and-Playable-Worlds-P01/main.pdf\#nameddest=appendix-49933872bb51}{Norms lengths and induced distances (Vision, XR, and Playable Worlds, Part 1, Reference notes: Norms lengths and induced distances)} --- Lengths and their domains are needed before unit directions or angular measurements.
\end{itemize}

\subsection{Used by}\label{appendix-8bc4e8fc547c-used-by}

\hypertarget{appendix-8bc4e8fc547c-used-by}{}

\begin{itemize}
\item
  \href{../02-Mathematics-for-Visual-Computing-P01/main.pdf\#nameddest=appendix-965d9260eeab}{Quaternion algebra and Hamilton product (Mathematics for Visual Computing, Part 1, Reference notes: Quaternion algebra and Hamilton product)}
\item
  \href{../07-Vision-XR-and-Playable-Worlds-P01/main.pdf\#nameddest=appendix-505a502c372a}{Projection rejection and scalar components (Vision, XR, and Playable Worlds, Part 1, Reference notes: Projection rejection and scalar components)} --- Builds on this mechanism in the connected geometry story.
\item
  \href{../07-Vision-XR-and-Playable-Worlds-P01/main.pdf\#nameddest=appendix-3b863809f82a}{Gram matrices and coordinate metrics (Vision, XR, and Playable Worlds, Part 1, Reference notes: Gram matrices and coordinate metrics)} --- Builds on this mechanism in the connected geometry story.
\item
  \href{../07-Vision-XR-and-Playable-Worlds-P01/main.pdf\#nameddest=appendix-2391db8052b1}{Cross products and oriented area (Vision, XR, and Playable Worlds, Part 1, Reference notes: Cross products and oriented area)} --- Builds on this mechanism in the connected geometry story.
\item
  \href{../07-Vision-XR-and-Playable-Worlds-P01/main.pdf\#nameddest=appendix-5caa058bd9c7}{Triple products and signed volume (Vision, XR, and Playable Worlds, Part 1, Reference notes: Triple products and signed volume)} --- Builds on this mechanism in the connected geometry story.
\item
  \href{../07-Vision-XR-and-Playable-Worlds-P01/main.pdf\#nameddest=appendix-215ac7099a82}{Minkowski sums and support mappings (Vision, XR, and Playable Worlds, Part 1, Reference notes: Minkowski sums and support mappings)} --- Builds on this mechanism in the connected geometry story.
\item
  \href{../07-Vision-XR-and-Playable-Worlds-P01/main.pdf\#nameddest=appendix-46da295c6787}{Covectors linear functionals and normal transforms (Vision, XR, and Playable Worlds, Part 1, Reference notes: Covectors linear functionals and normal transforms)} --- Builds on this mechanism in the connected geometry story.
\item
  \href{../07-Vision-XR-and-Playable-Worlds-P01/main.pdf\#nameddest=appendix-c3421e6b4f2c}{QR factorization and least squares (Vision, XR, and Playable Worlds, Part 1, Reference notes: QR factorization and least squares)} --- Extends this foundation into linear-systems and decomposition analysis.
\item
  \href{../04-Light-Transport-and-Image-Synthesis-P01/main.pdf\#nameddest=appendix-1fad5e7f36fa}{Singular value decomposition (Light Transport and Image Synthesis, Part 1, Reference notes: Singular value decomposition)} --- Extends this foundation into linear-systems and decomposition analysis. \#\# Recall and next depth
\end{itemize}

Derive the cosine relation and Cauchy--Schwarz; use a scale-aware orthogonality criterion. Full production proofs, target-specific integration and learner mastery require further evidence.

\subsection{Source and navigation}\label{appendix-8bc4e8fc547c-source-and-navigation}

\hypertarget{appendix-8bc4e8fc547c-source-and-navigation}{}

\href{../02-Mathematics-for-Visual-Computing-P01/main.pdf\#nameddest=01-foundations-i-ch017-the-dot-product-turns-two-movements-into-a-measurement}{Teaching narrative (Mathematics for Visual Computing, Part 1, Ch. 1)} · \hyperref[appendix-fd090d982a15-phase-6---parts-501-600]{Coverage and source gap} · \hyperref[appendix-a103e7846d21-phase-6---source-analysis-and-verification]{Evidence and limits} · \href{../07-Vision-XR-and-Playable-Worlds-P01/main.pdf\#nameddest=appendix-4ad47bdd269a}{Dependency Map (Vision, XR, and Playable Worlds, Part 1, Reference notes: Dependency Map)}
